\chapter{Какво беше необходимо, за да работят технологиите}

Една от важните части на заданието е да се опише не само видимата функционалност, но и интеграционната работа, без която крайният резултат не би бил възможен. В настоящия проект именно тази невидима на пръв поглед работа се оказа съществена. Архитектурно кодът можеше да изглежда правилен още в ранна фаза, но реалното съгласуване между уеб и настолния клиент, между двете бази данни и между обработката на видео и потребителския интерфейс изискваше допълнителни решения.

\section{Съгласуване на конфигурацията}

За да могат двата клиента да показват еднакви данни и да работят последователно, се оказа необходимо и двата да сочат към една и съща конфигурационна среда. Това включва не само избора на \texttt{Persistence:Provider}, но и еднакъв низ за свързване, еднакво файлово хранилище, еднакъв път към ONNX модела и еднаква конфигурация за \texttt{ffmpeg}. Без това приложението би изглеждало функционално, но двата клиента биха могли да показват различни данни и така да създадат погрешно впечатление за архитектурно разминаване.

Тази особеност е важна и от гледна точка на защитата. Изискването за две UI технологии не означава две независими програми, а две визуализации върху споделена приложна логика. Аналогично, изискването за две СУБД не означава два отделни проекта, а един модел и един бизнес слой, който работи и с двата доставчика. Следователно правилната конфигурация е част от архитектурната демонстрация, а не просто техническа подробност.

\section{Стартиране на настолния клиент и фонови услуги}

При Avalonia се появи специфичен проблем. Ако настолното приложение използва само ръчно конструиран \texttt{ServiceCollection}, фонова услуга като обработчика на анализи няма да бъде стартирана по същия начин, както в уеб приложението. За да се избегне това разминаване, Avalonia клиентът беше свързан към \texttt{IHost}. По този начин настолното приложение започна да използва същия модел за инжектиране на зависимости, фонови услуги и конфигурация, който вече съществуваше в уеб слоя.

Това решение се оказа ключово, защото позволи \texttt{QueuedAnalysisBackgroundService} да работи и в настолната среда и доближи стартирането на двата клиента. Така Avalonia приложението се превърна в реално функционален интерфейс, а не само във формално доказателство за втора технология.

\section{Медийна последователност и обработка на видео}

Друг важен слой от работа беше свързан с реалната обработка на видео. За да не остане приложението на ниво административен интерфейс с фиктивни резултати, беше необходимо качените файлове действително да се записват във файлово хранилище, от тях да се извличат кадри чрез \texttt{ffmpeg}, тези кадри да се подават към YOLO модела и накрая да се генерира анотирано изходно видео.

Тази интеграция не е видима директно от структурата на решението, но е сърцевината на функционалността. Тя обвързва файловата система, външния инструмент за обработка на видео, модула за инференция и слоя за съхранение в един последователен процес. Без такава реализация приложението би имало добра архитектура, но не би показвало реален аналитичен резултат.

\section{Работа под Linux и локална среда}

Фактът, че разработката и тестването се извършват под Linux, постави и няколко допълнителни изисквания. Наложи се да бъде конфигуриран коректен път към \texttt{/usr/bin/ffmpeg}, да се вземе предвид поведението на локалните HTTPS сертификати в браузъра и да се намери подходящ настолен подход за визуализация на видео без въвеждане на тежка зависимост за вградено възпроизвеждане. В резултат настолният клиент комбинира предварителни кадри и отваряне на оригиналното или обработеното видео чрез системния видео плейър.

Тук ясно се вижда разликата между това кодът да е написан и това приложението да е действително работещо в реална среда. Именно тези интеграционни аспекти често не се забелязват в крайния продукт, но представляват съществена част от извършената инженерна работа.
